\emph{Preliminaries.} Everything about the adjoint, the scalar part and states is taken from Theorem \claimref{synthesis.chsh-bound-for-two-subsystems-of-a-clifford-algebra-by-a-rational-certificate} (the base): the blades are a basis (Theorem \claimref{synthesis.signed-directions-and-the-integer-clifford-algebra}, with rational coefficients by Theorem \claimref{synthesis.division-gives-the-rationals}), $\langle x^{\dagger}x\rangle_{0}$ is the sum of the squares of the coefficients, $\langle xy\rangle_{0} = \langle yx\rangle_{0}$, a functional $\langle\rho\,\cdot\,\rangle_{0}/\langle\rho\rangle_{0}$ with $\rho = \rho\rho^{\dagger}$ is a state, and $\omega(h)^{2} \le \omega(h^{2})$ for a state and self-adjoint $h$. By clause (vii) of the base, $A_{1}, A_{2}, B_{1}, B_{2}$ are observables, anticommuting within the pairs and commuting across; by its clauses (iv) and (viii), $C^{2} = 4 + 4K$ with $K = -A_{1}A_{2}B_{1}B_{2} = e_{1}f_{2}$, $K^{\dagger} = K$, $K^{2} = 1$, $KC = CK = C$, $K$ commutes with each $A_{i}B_{j}$, and $\langle C\rangle_{0} = \langle K\rangle_{0} = 0$.

\emph{(i).} $I^{2} = B_{1}B_{2}B_{1}B_{2} = -1$, $I^{\dagger} = B_{2}B_{1} = -I$; $I$ commutes with the $A$'s, anticommutes with $B_{1}$ and $B_{2}$, and commutes with $K = -A_{1}A_{2}I$. $B_{1}I = B_{2}$ and $B_{2}I = -B_{1}$, so $CI = A_{1}(B_{2} - B_{1}) + A_{2}(B_{2} + B_{1})$, which is the stated form and is the CHSH element $A'_{1}(B_{1} + B_{2}) + A'_{2}(B_{1} - B_{2})$ for $A'_{1} = A_{2}$, $A'_{2} = -A_{1}$. Each term of $C$ contains one $B_{j}$, so $IC = -CI$; then $(C^{\perp})^{\dagger} = I^{\dagger}C = -IC = CI$, $CC^{\perp} = C^{2}I$ and $C^{\perp}C = CIC = -C^{2}I$, and $(C^{\perp})^{2} = CICI = -C^{2}I^{2} = C^{2}$. Hence $(xC + yC^{\perp})^{2} = (x^{2} + y^{2})C^{2} = 4(x^{2} + y^{2})(1 + K)$. For a state put $X = \omega(C)$, $Y = \omega(C^{\perp})$ and $h = XC + YC^{\perp}$: $\omega(h)^{2} \le \omega(h^{2})$ reads $(X^{2} + Y^{2})^{2} \le 4(X^{2} + Y^{2})(1 + \omega(K))$, and $\omega(K)^{2} \le \omega(K^{2}) = 1$; so $X^{2} + Y^{2} \le 8$.

\emph{(ii).} \emph{The corner.} For a blade $b$, $KbK = \pm b$, so $P_{K}bP_{K}$ is $bP_{K}$ if $b$ commutes with $K$ and $0$ otherwise; half of the sixteen blades commute with the blade $K$, and $b$ and $bK$ give the same element up to sign, so the corner has dimension at most $4$. $P_{K}$, $C = CP_{K}$, $IP_{K}$, $C^{\perp} = C^{\perp}P_{K}$ lie in it ($KC = C$, and $K$ commutes with $I$) and are linearly independent (the check computes the rank $4$; alternatively they are a unit, two anticommuting elements of squares $8P_{K}$ and $-P_{K}$, and their product). \emph{Idempotents.} The self-adjoint elements of the corner are $t = \sigma P_{K} + xC + yC^{\perp}$, since $IP_{K}$ is anti-self-adjoint. By (i), $t^{2} = (\sigma^{2} + 8(x^{2} + y^{2}))P_{K} + 2\sigma(xC + yC^{\perp})$, using $(1 + K)P_{K} = 2P_{K}$. So $t^{2} = t$ if and only if $x = y = 0$ and $\sigma \in \{0, 1\}$, or $\sigma = \tfrac12$ and $32(x^{2} + y^{2}) = 1$; with $X = 16x$, $Y = 16y$ this is $X^{2} + Y^{2} = 8$. \emph{Values.} $\langle P_{K}\rangle_{0} = \tfrac12$ and $\langle C\rangle_{0} = \langle C^{\perp}\rangle_{0} = \langle CC^{\perp}\rangle_{0} = 0$ ($C$ and $C^{\perp}$ are sums of blades other than $1$, and $\langle CC^{\perp}\rangle_{0} = \langle C^{\perp}C\rangle_{0} = -\langle CC^{\perp}\rangle_{0}$ by cyclicity and anticommutation), and $\langle C^{2}\rangle_{0} = 4$. So $\langle\rho_{X,Y}\rangle_{0} = \tfrac14$, $4\langle\rho_{X,Y}C\rangle_{0} = 4 \cdot \tfrac{X}{16} \cdot 4 = X$ (the term $\tfrac12P_{K}C = \tfrac12C$ has no scalar part) and likewise $4\langle\rho_{X,Y}C^{\perp}\rangle_{0} = Y$. \emph{Purity.} Let $\rho = \rho_{X,Y}$ and $\rho' = P_{K} - \rho = \rho_{-X,-Y}$, so $\rho\rho' = 0$. Since $I$ anticommutes with $C$ and $C^{\perp}$, $I\rho' = \rho I$; hence $\rho I\rho' = \rho I \ne 0$ ($I$ is invertible) and likewise $\rho'I\rho = \rho'I \ne 0$. The corner is the direct sum of $\rho\,\cdot\,\rho$, $\rho\,\cdot\,\rho'$, $\rho'\,\cdot\,\rho$, $\rho'\,\cdot\,\rho'$ applied to it, each of which is nonzero (they contain $\rho$, $\rho I$, $\rho'I$, $\rho'$), so each has dimension $1$. Therefore $\rho x\rho = \rho(P_{K}xP_{K})\rho = \lambda(x)\rho$, and taking scalar parts with cyclicity, $\langle\rho x\rangle_{0} = \lambda(x)\langle\rho\rangle_{0} = \lambda(x)/4$. $\omega_{\rho}$ is a state since $\rho = \rho\rho^{\dagger}$. Conversely a pure state with $K\rho = \rho$ has $\rho = \rho^{\dagger}$, so $\rho K = \rho$ and $\rho$ is a self-adjoint idempotent of the corner with $\langle\rho\rangle_{0} = \tfrac14$, which excludes $0$ and $P_{K}$. A rational point with $Y = 0$ would have $X^{2} = 8$.

\emph{(iii).} $\rho_{0} = \tfrac14(1 + A_{1}B_{1} - A_{2}B_{2} - A_{1}B_{1}A_{2}B_{2})$ and $A_{1}B_{1}A_{2}B_{2} = A_{1}A_{2}B_{1}B_{2} = -K$; $C - C^{\perp} = 2(A_{1}B_{1} - A_{2}B_{2})$. So $\rho_{0} = \tfrac12P_{K} + \tfrac18(C - C^{\perp}) = \rho_{2,-2}$. $y_{u}^{\dagger} = 1 - uI$ and $y_{u}^{\dagger}y_{u} = 1 + u^{2}$, so $y_{u}^{-1} = y_{u}^{\dagger}/(1 + u^{2})$ and $\rho_{u} = y_{u}\rho_{0}y_{u}^{-1}$. $y_{u}$ commutes with $P_{K}$, and $y_{u}C = Cy_{u}^{\dagger}$ since $IC = -CI$; hence $y_{u}Cy_{u}^{-1} = Cy_{u}^{\dagger}y_{u}^{-1} = CR_{u}$ with $R_{u} = (1 - uI)(1 + uI)^{-1} = c_{u} - s_{u}I$, the formula of Theorem \claimref{synthesis.the-sign-is-the-rotor-at-infinity}, which uses only $I^{2} = -1$. So $y_{u}Cy_{u}^{-1} = c_{u}C - s_{u}C^{\perp}$ and $y_{u}C^{\perp}y_{u}^{-1} = CR_{u}I = c_{u}C^{\perp} + s_{u}C$. Then $\rho_{u} = \tfrac12P_{K} + \tfrac18\bigl((c_{u} - s_{u})C - (c_{u} + s_{u})C^{\perp}\bigr) = \rho_{X,Y}$ with $(X, Y) = (2(c_{u} - s_{u}), -2(c_{u} + s_{u}))$, and $2(c_{u} - s_{u}) = 2(1 - 2u - u^{2})/(1 + u^{2})$. For $u = \infty$: $ICI^{-1} = -C$, $IC^{\perp}I^{-1} = -C^{\perp}$, so $I\rho_{0}I^{\dagger} = \rho_{-2,2}$, the limit $(c_{u}, s_{u}) = (-1, 0)$. The map $(c', s') \mapsto (2(c' - s'), -2(c' + s'))$ is a bijection from the rational points of the unit circle onto those of $X^{2} + Y^{2} = 8$ (inverse $c' = (X - Y)/4$, $s' = -(X + Y)/4$), and $u \mapsto (c_{u}, s_{u})$ is a bijection from $\Q \cup \{\infty\}$ onto the rational points of the unit circle by clause (iii) of that Theorem; so $u \mapsto \rho_{u}$ is a bijection onto the pure states of (ii). \emph{Turned directions.} By cyclicity $\omega_{\rho_{u}}(C) = 4\langle\rho_{0}\,y_{u}^{-1}Cy_{u}\rangle_{0}$, and $y_{u}$ commutes with the $A$'s, so $y_{u}^{-1}Cy_{u}$ is the CHSH element of $A_{1}, A_{2}, b_{1}, b_{2}$ with $b_{j} = y_{u}^{-1}B_{j}y_{u} = B_{j}(y_{u}^{\dagger})^{-1}y_{u} = B_{j}R_{u}^{-1}$ ($I$ anticommutes with $B_{j}$) and $R_{u}^{-1} = c_{u} + s_{u}I$: $b_{1} = c_{u}B_{1} + s_{u}B_{2}$, $b_{2} = c_{u}B_{2} - s_{u}B_{1}$, self-adjoint of square one, commuting with the $A$'s. \emph{Parity.} $\Gamma = p_{1}p_{2}$ anticommutes with $A_{1} = e_{1}$ and $B_{1} = e_{2}p_{1}$ and commutes with $A_{2}$, $B_{2}$, $K$. So $\Gamma C\Gamma = A_{1}B_{1} - A_{1}B_{2} - A_{2}B_{1} - A_{2}B_{2} = -C^{\perp}$ and $\Gamma C^{\perp}\Gamma = -C$, whence $\Gamma\rho_{X,Y}\Gamma = \rho_{-Y,-X}$, which equals $\rho_{X,Y}$ exactly when $Y = -X$, that is $X = \pm 2$. With $v_{k} = \tfrac12(e_{k} + f_{k})$, $w_{k} = \tfrac12(e_{k} - f_{k})$ the identity $\tfrac14(1 + p_{1}p_{2}) - \tfrac12(w_{1}w_{2} + v_{2}v_{1}) = \rho_{\infty}$ is an identity between explicit elements, verified by the check.

\emph{(iv).} If $P$ is odd, $P^{2} \equiv 1 \pmod 8$ and $8Q^{2} - P^{2} \equiv 7 \pmod 8$, so it is at least $7$. If $P = 2m$, then $8Q^{2} - P^{2} = 4(2Q^{2} - m^{2})$ is a positive multiple of $4$, equal to $4$ exactly when $m^{2} - 2Q^{2} = -1$. A state value is rational and by clause (iii) of the base its square is below $8$; write it $\pm P/Q$.

\emph{(v).} If $x^{2} - 2Q^{2} = -1$ then $(3x + 4Q)^{2} - 2(2x + 3Q)^{2} = x^{2} - 2Q^{2}$, so the recursion produces solutions. Conversely let $(x, Q)$ be a positive solution with $Q > 1$; then $Q \ge 3$ ($Q = 2$ gives $x^{2} = 7$), and $x_{-} = 3x - 4Q$, $Q_{-} = 3Q - 2x$ is a solution with $(3x_{-} + 4Q_{-}, 2x_{-} + 3Q_{-}) = (x, Q)$. From $x^{2} = 2Q^{2} - 1$: $x > Q$ gives $Q_{-} < Q$; $4x^{2} < 9Q^{2}$ gives $Q_{-} > 0$; and $9x^{2} = 18Q^{2} - 9 > 16Q^{2}$ for $Q \ge 3$ gives $x_{-} > 0$. By descent every positive solution is in the sequence. For a solution, $x$ is odd, $a = (x - 1)/2$ and $b = (x + 1)/2$ satisfy $a^{2} + b^{2} = (x^{2} + 1)/2 = Q^{2}$, so $(c_{u}, s_{u}) = (b/Q, -a/Q)$ is a rational point of the unit circle with $2(c_{u} - s_{u}) = 2x/Q$ and $s_{u} < 0$: by (iii) $S = 2x/Q$ is the value of a pure state, with $Y = -2(c_{u} + s_{u}) = -2/Q$ (and $(a/Q, -b/Q)$ gives $Y = 2/Q$). $8 - S^{2} = 4(2Q^{2} - x^{2})/Q^{2} = 4/Q^{2}$, and $(S + Q^{-2})^{2} = 8 - 4Q^{-2} + 2SQ^{-2} + Q^{-4} > 8$ because $S \ge 2$. If a state value has denominator $Q' < Q_{k}$, its squared deficit is at least $4/Q'^{2} > 4/Q_{k}^{2}$ by (iv), so its absolute value is below $S_{k}$. $Q_{k+1} > Q_{k}$, so $S_{k}$ increases and $8 - S_{k}^{2} \to 0$; with clause (iii) of the base, the supremum of the values at these settings is $2\sqrt{2}$ and is not attained. $\tfrac{280}{99}$ is a value by clause (iv) of the base ($r = \tfrac{35}{99}$, $8r^{2} = \tfrac{9800}{9801}$), and $280 \cdot 29 = 8120 > 8118 = 82 \cdot 99$. The last sentence of the clause is an enumeration by the check: for each denominator up to $400$, the fractions $P/Q$ in lowest terms with $8Q^{2} - P^{2}$ a perfect square (the values $X$ of rational points $(P/Q, Y)$ of the circle), and their running maximum.
